\section{非线性 RL 广义单调性与 PPA 的局部收敛}
\label{sec:deterministic-rleb-spine}

设 $X=\mathbb R^n$，$F:X\rightrightarrows X$，$\lambda>0$，
并取非空闭集 $S\subset F^{-1}(0)$。沿用引言中
$J_{\lambda F}$ 的记号，以下记
\[
 d(x,A):=\inf_{a\in A}\|x-a\|,
 \qquad
 r_F(y):=d(0,F(y)),
 \qquad
 P_S(x):=\{p\in S:\|x-p\|=d(x,S)\},
\]
并约定 $d(x,\varnothing)=+\infty$。由于 $S$ 非空且闭，
$P_S(x)$ 非空，但未必单值。

\subsection{RL 条件及其 \texorpdfstring{$J$}{J} 侧形式}

\begin{definition}[RL 广义单调性]
\label{def:graph-RL}
设 $\mathcal G\subset\operatorname{gph}F$，$R>0$，$L\geq0$，
$0<\gamma\leq1$。称 $\mathcal G$ 在尺度 $R$ 上满足
$\mathrm{RL}(\lambda,\gamma,L)$，如果对任意
$(u,u^*),(v,v^*)\in\mathcal G$，记
\[
 a:=u-v,\qquad b:=u^*-v^*,
\]
只要 $\|a+\lambda b\|\leq R$，就有
\begin{equation}
\label{eq:graph-RL}
 \lambda\langle a,b\rangle
 \geq\frac14\left(
   \|a+\lambda b\|^2-L^2\|a+\lambda b\|^{2\gamma}
 \right).
\end{equation}
由恒等式
$\|a+\lambda b\|^2-\|a-\lambda b\|^2
 =4\lambda\langle a,b\rangle$，
这一条件等价于
\begin{equation}
\label{eq:power-RL}
 \|a-\lambda b\|\leq L\|a+\lambda b\|^\gamma
\end{equation}
对上述点对成立。
\end{definition}

固定 $\lambda$ 时，简记
$\mathrm{RL}(\lambda,\gamma,L)$ 为 $\mathrm{RL}(\gamma,L)$。
所检验的点对可以来自不同分支；当 $\gamma=L=1$ 时，
条件成为这些点对上的单调性不等式。

与图像局部化 $\mathcal G$ 对应的近端关系定义为
\[
 J_{\mathcal G}(x)
 :=\{u:(u,(x-u)/\lambda)\in\mathcal G\}.
\]
恒有 $J_{\mathcal G}(x)\subset J_{\lambda F}(x)$。
若某输入区域内所有近端输出对应的图像点均属于
$\mathcal G$，则两者在该区域上一致。

\begin{lemma}[RL 条件的 $J$ 侧形式]
\label{lem:minty-equivalence}
定义~\ref{def:graph-RL} 中的 RL 条件等价于：对任意
$u\in J_{\mathcal G}(x)$ 和 $v\in J_{\mathcal G}(y)$，
只要 $\|x-y\|\leq R$，就有
\begin{equation}
\label{eq:resolvent-energy-all-pairs}
 \|u-v\|^2+\|(x-u)-(y-v)\|^2
 \leq\frac12\left(
   \|x-y\|^2+L^2\|x-y\|^{2\gamma}
 \right).
\end{equation}
特别地，$J_{\mathcal G}$ 在其定义域上单值。
\end{lemma}

\begin{proof}
对相应的图像点，记
$a=u-v$，$\lambda b=(x-u)-(y-v)$，则
$a+\lambda b=x-y$，且
\[
 \|a\|^2+\lambda^2\|b\|^2
 =\|a+\lambda b\|^2-2\lambda\langle a,b\rangle.
\]
将此式代入 \eqref{eq:graph-RL}，即得
\eqref{eq:resolvent-energy-all-pairs}；反向代入亦成立。
令 $x=y$，上式给出 $2\|u-v\|^2\leq0$，故 $u=v$。
\end{proof}

\begin{definition}[广义度量次正则性]
\label{def:admissible-gauge}
设 $V\subset X$，$\bar t>0$，且
$\psi:[0,\bar t]\to[0,\infty)$ 单调不减，
满足 $\psi(0)=0$ 及 $\psi(t)\to0$（$t\downarrow0$）。
若
\begin{equation}
\label{eq:gauge-EB}
 d(y,S)\leq\psi\bigl(r_F(y)\bigr)
 \qquad (y\in V,\ r_F(y)\leq\bar t),
\end{equation}
则称 $F$ 在 $V$ 上满足相对于 $S$ 的
$\psi$-广义度量次正则性。
\end{definition}

\begin{assumption}[局部 RLEB 假设]
\label{ass:local-rleb}
设 $U\subset X$ 为开集，$\mathcal G\subset\operatorname{gph}F$，
并取
\[
R>0,\quad \bar t>0,\quad L\geq0,\quad
0<\gamma\leq1,\quad 0<\kappa<1,
\]
以及定义~\ref{def:admissible-gauge} 中的控制函数 $\psi$。
记 $U_R:=\{x\in U:d(x,S)\leq R\}$。
\end{assumption}
对 $(u,u^*)\in\mathcal G$，记
\[
 M_+(u,u^*):=u+\lambda u^*,
 \qquad
 M_+(\mathcal G)
 :=\{M_+(u,u^*):(u,u^*)\in\mathcal G\}.
\]
由 $J_{\mathcal G}$ 的定义可知，
$M_+(\mathcal G)=\operatorname{dom}J_{\mathcal G}$。
以下是局部近端迭代所需的条件。

\begin{enumerate}
\item[(A1)] 
$U_R\subset M_+(\mathcal G)$；并且对任意
$x\in U_R$，存在 $p\in P_S(x)$ 使得
$(p,0)\in\mathcal G$。

\item[(A2)] 
$\mathcal G$ 在尺度 $R$ 上满足
$\mathrm{RL}(\lambda,\gamma,L)$ 条件。

\item[(A3)] 
对任意 $y\in J_{\mathcal G}(U_R)$，
若 $r_F(y)\leq\bar t$，则
$d(y,S)\leq\psi\bigl(r_F(y)\bigr)$；
此外，
\begin{equation}
\label{eq:gauge-domain}
 \frac{R+LR^\gamma}{2\lambda}\leq\bar t.
\end{equation}

\item[(A4)] 
对任意 $0<r\leq R$，有
\begin{equation}
\label{eq:direct-compatibility}
 \psi\!\left(\frac{r+Lr^\gamma}{2\lambda}\right)
 \leq\kappa r.
\end{equation}
\end{enumerate}

条件 (A1) 中的值域覆盖要求保证
$J_{\mathcal G}$ 在 $U_R$ 上处处有定义；
类似的局部覆盖条件见
\cite[Assumption~2(d)(ii)]{LukeTam2025GeneralizedMonotonicity}。
(A1) 中关于最近零点的要求用于下文的一步比较。
根据引理~\ref{lem:minty-equivalence}，(A2) 保证
$J_{\mathcal G}$ 在其定义域上单值。因此，
(A1)--(A2) 确定了 $U_R$ 上的单值局部近端映射；
(A3)--(A4) 分别用于输出点的误差估计和距离递推的收缩。

\begin{lemma}[RLEB 的一步估计]
\label{lem:solution-comparison}
设 \emph{(A1)--(A3)} 成立。取 $x\in U_R$，
令 $x^+$ 为 $J_{\mathcal G}(x)$ 的唯一元素，并记
\[
 d:=d(x,S),\qquad
 d^+:=d(x^+,S),\qquad
 s:=\|x-x^+\|.
\]
则
\begin{align}
\label{eq:energy-one-step}
 (d^+)^2+s^2
 &\leq\frac{d^2+L^2d^{2\gamma}}2,\\
\label{eq:step-one-step}
 s&\leq\frac{d+Ld^\gamma}{2},\\
\label{eq:direct-output-bound}
 d^+&\leq\psi(s/\lambda),\\
\label{eq:direct-recursion-general}
 d^+&\leq
 \psi\!\left(\frac{d+Ld^\gamma}{2\lambda}\right).
\end{align}
\end{lemma}

\begin{proof}
由 (A1) 取 $p\in P_S(x)$，使得 $(p,0)\in\mathcal G$。
由于 $\|x-p\|=d\leq R$，将
$(x^+,(x-x^+)/\lambda)$ 与 $(p,0)$
代入引理~\ref{lem:minty-equivalence}，得到
\[
 \|x^+-p\|^2+s^2
 \leq\frac{d^2+L^2d^{2\gamma}}2.
\]
结合 $d^+\leq\|x^+-p\|$，
即得 \eqref{eq:energy-one-step}。

再由 $x^+-p=(x-p)-(x-x^+)$
及 Cauchy--Schwarz 不等式，有
\[
 d^2-2ds+2s^2
 \leq\|x^+-p\|^2+s^2
 \leq\frac{d^2+L^2d^{2\gamma}}2.
\]
整理得
\begin{equation}
\label{eq:power-step}
 (2s-d)^2\leq L^2d^{2\gamma},
 \qquad
 s\leq\frac{d+Ld^\gamma}{2},
\end{equation}
从而证明 \eqref{eq:step-one-step}。

由 $(x-x^+)/\lambda\in F(x^+)$ 可知
$r_F(x^+)\leq s/\lambda$。同时，
\eqref{eq:step-one-step} 和
\eqref{eq:gauge-domain} 给出
\[
 \frac{s}{\lambda}
 \leq\frac{d+Ld^\gamma}{2\lambda}
 \leq\frac{R+LR^\gamma}{2\lambda}
 \leq\bar t.
\]
因此，(A3) 中的误差界适用于 $x^+$。
利用 $\psi$ 的单调性，得到
\[
 d^+\leq\psi\bigl(r_F(x^+)\bigr)
 \leq\psi(s/\lambda)
 \leq
 \psi\!\left(\frac{d+Ld^\gamma}{2\lambda}\right),
\]
即 \eqref{eq:direct-output-bound} 和
\eqref{eq:direct-recursion-general}。
\end{proof}

\subsection{局部收敛定理}

\begin{theorem}[非线性 RL 与广义误差界下的局部收敛]
\label{thm:two-branch-RLEB}
设假设~\ref{ass:local-rleb} 成立。对 $r\geq0$，记
\[
 \mathcal L(r):=
 \frac12\left(
   \frac{r}{1-\kappa}
   +\frac{Lr^\gamma}{1-\kappa^\gamma}
 \right).
\]
取 $x^0\in U_R$，令 $d_0:=d(x^0,S)$。若
\begin{equation}
\label{eq:direct-localization}
 \mathcal L(d_0)<d(x^0,X\setminus U),
\end{equation}
则局部近端迭代
\begin{equation}
\label{eq:PPA}
 x^{k+1}=J_{\mathcal G}(x^k)
\end{equation}
可无限进行，所有迭代点均属于 $U_R$，且轨道具有有限长度，
收敛到某个 $x^\infty\in S$。对任意 $k\geq0$，还有
\begin{align}
\label{eq:direct-distance-rate}
 d(x^k,S)&\leq\kappa^k d_0,\\
\label{eq:direct-step-rate}
 \|x^{k+1}-x^k\|
 &\leq\frac12\bigl(\kappa^k d_0
                   +L(\kappa^k d_0)^\gamma\bigr),\\
\label{eq:direct-tail}
 \|x^\infty-x^k\|
 &\leq\mathcal L(\kappa^k d_0).
\end{align}
\end{theorem}

\begin{proof}
设 $x^k\in U_R$，记
$d_k:=d(x^k,S)$，
$s_k:=\|x^{k+1}-x^k\|$。
由 (A1)--(A2)，$x^{k+1}=J_{\mathcal G}(x^k)$
存在且唯一。若 $d_k>0$，则引理~\ref{lem:solution-comparison}
与 (A4) 给出
\[
 d_{k+1}
 \leq\psi\!\left(
   \frac{d_k+Ld_k^\gamma}{2\lambda}
 \right)
 \leq\kappa d_k.
\]
若 $d_k=0$，则由 \eqref{eq:step-one-step} 得
$s_k=0$，故仍有 $d_{k+1}\leq\kappa d_k$。

因此，对所有已定义的迭代点，
\[
 d_k\leq\kappa^k d_0,\qquad
 s_k\leq\frac12\bigl(
   \kappa^k d_0+L\kappa^{\gamma k}d_0^\gamma
 \bigr).
\]
特别地，对任意 $n\geq0$，
\[
 \|x^{n+1}-x^0\|
 \leq\sum_{k=0}^{n}s_k
 \leq\mathcal L(d_0)
 <d(x^0,X\setminus U).
\]
于是 $x^{n+1}\in U$；又因
$d_{n+1}\leq\kappa d_n\leq R$，
故 $x^{n+1}\in U_R$。归纳表明迭代可无限进行，
并证明 \eqref{eq:direct-distance-rate}
和 \eqref{eq:direct-step-rate}。

最后，对任意 $k\geq0$，有
\[
 \sum_{j=k}^{\infty}s_j
 \leq\mathcal L(\kappa^k d_0).
\]
因此轨道具有有限长度并收敛到某个 $x^\infty\in X$，
且满足 \eqref{eq:direct-tail}。
由于 $d(x^k,S)\to0$ 且 $S$ 为闭集，
可得 $x^\infty\in S$。
\end{proof}

\begin{corollary}[集值近端关系的推广]
\label{cor:selection-extension}
沿用假设~\ref{ass:local-rleb} 中的集合、参数及控制函数，
并设 $\mathcal J(x)\subset J_{\lambda F}(x)$
为 $U_R$ 上的近端关系。以下列条件替代
\emph{(A1)--(A3)}：
\begin{enumerate}
\item 对任意 $x\in U_R$，$\mathcal J(x)\ne\varnothing$；

\item 对任意 $x\in U_R$ 及
$x^+\in\mathcal J(x)$，存在 $p\in P_S(x)$
使得
\begin{equation}
\label{eq:solution-comparison-innerproduct}
 \langle x^+-p,x-x^+\rangle
 \geq\frac14\left(
   d(x,S)^2-L^2d(x,S)^{2\gamma}
 \right);
\end{equation}

\item 对任意 $x\in U_R$ 及
$x^+\in\mathcal J(x)$，若
$r_F(x^+)\leq\bar t$，则
$d(x^+,S)\leq\psi\bigl(r_F(x^+)\bigr)$；
同时，\eqref{eq:gauge-domain} 成立。
\end{enumerate}

若 \emph{(A4)} 成立，且 $x^0\in U_R$ 满足
局部化条件 \eqref{eq:direct-localization}，则任意按
$x^{k+1}\in\mathcal J(x^k)$ 选取的轨道均可无限延拓，
始终位于 $U_R$，具有有限长度，并收敛到 $S$ 中某一点。
其距离、步长及尾项分别满足
\eqref{eq:direct-distance-rate}--\eqref{eq:direct-tail}。
这些估计对所有选取采用相同的常数；极限点可以随选取而改变。
\end{corollary}

\begin{proof}
任取 $x\in U_R$ 和 $x^+\in\mathcal J(x)$，
按条件取 $p\in P_S(x)$，并记
\[
 d:=d(x,S),\qquad
 d^+:=d(x^+,S),\qquad
 s:=\|x-x^+\|.
\]
由 $x-p=(x^+-p)+(x-x^+)$ 及
\eqref{eq:solution-comparison-innerproduct}，有
\[
 \|x^+-p\|^2+s^2
 =d^2-2\langle x^+-p,x-x^+\rangle
 \leq\frac{d^2+L^2d^{2\gamma}}2.
\]
另一方面，Cauchy--Schwarz 不等式给出
\[
 \frac14\bigl(d^2-L^2d^{2\gamma}\bigr)
 \leq\langle x^+-p,x-x^+\rangle
 \leq ds-s^2.
\]
因此
\[
 (2s-d)^2\leq L^2d^{2\gamma},
 \qquad
 s\leq\frac{d+Ld^\gamma}{2}.
\]
由 $x^+\in J_{\lambda F}(x)$ 可知
$r_F(x^+)\leq s/\lambda$。结合条件中的
\eqref{eq:gauge-domain} 与误差界，得到
\[
 d^+\leq
 \psi\!\left(\frac{d+Ld^\gamma}{2\lambda}\right).
\]
上述估计对每个允许的 $x^+$ 均成立。
利用 (A4)、各点上 $\mathcal J$ 的非空性及
\eqref{eq:direct-localization}，按定理~
\ref{thm:two-branch-RLEB} 证明中的归纳与求和步骤，
即可得到所述结论。
\end{proof}

若允许缩小工作半径 $R$ 并重新选取 $\kappa$，
则 (A4) 在某个充分小的半径上成立，当且仅当
\begin{equation}
\label{eq:direct-limsup}
 \limsup_{r\downarrow0}
 \frac{\psi\bigl((r+Lr^\gamma)/(2\lambda)\bigr)}{r}<1.
\end{equation}
事实上，若上式成立，可在该极限上确界与 $1$ 之间
选取 $\kappa$，再取充分小的 $R$，使
\eqref{eq:direct-compatibility} 成立；反之，
(A4) 直接给出上式。由于
$(R+LR^\gamma)/(2\lambda)\to0$（$R\downarrow0$），
进一步缩小 $R$ 还可保证
\eqref{eq:gauge-domain}。

\begin{corollary}[残差增长条件与临界阶]
\label{cor:residual-growth}
设 \emph{(A1)--(A2)} 成立，且
$0<\gamma<1$、$L>0$。进一步假设存在
$m>0$ 和 $a>0$，使得对所有满足
$y\in J_{\mathcal G}(U_R)$、
$r_F(y)\leq\bar t$ 的点 $y$，有
\begin{equation}
\label{eq:residual-growth}
 r_F(y)\geq m\,d(y,S)^a.
\end{equation}
则在 $[0,\bar t]$ 上取
$\psi(t)=(t/m)^{1/a}$，可得到近端输出上的误差界。

对于这一控制函数，局部相容性条件
\eqref{eq:direct-limsup} 在 $a<\gamma$ 时成立；
在 $a=\gamma$ 时，成立当且仅当
\begin{equation}
\label{eq:H-critical-threshold}
 \lambda m>\frac L2;
\end{equation}
在 $a>\gamma$ 时则不成立。

在前两种相容情形下，适当缩小工作半径，
使 \eqref{eq:gauge-domain} 成立，并选取满足
\eqref{eq:direct-localization} 的初始点，
即可应用定理~\ref{thm:two-branch-RLEB}。
当 $a>\gamma$ 时，上述标量判据失效；
这本身不意味着近端迭代发散。
\end{corollary}

\begin{proof}
由 \eqref{eq:residual-growth} 可得，在所述输出区域上，
\[
 d(y,S)\leq
 \left(\frac{r_F(y)}{m}\right)^{1/a}.
\]
由于 $m,a>0$，函数
$\psi(t)=(t/m)^{1/a}$ 满足控制函数的要求。

对 $r>0$，\eqref{eq:direct-limsup} 中的商可写为
\[
 \frac{\psi\bigl((r+Lr^\gamma)/(2\lambda)\bigr)}{r}
 =
 \left(\frac{L+r^{1-\gamma}}{2\lambda m}\right)^{1/a}
 r^{\gamma/a-1}.
\]
因 $0<\gamma<1$ 且 $L>0$，右侧第一个因子
趋于正常数。故当 $a<\gamma$ 时，该商趋于 $0$；
当 $a>\gamma$ 时趋于 $+\infty$；当 $a=\gamma$ 时趋于
\[
 \left(\frac{L}{2\lambda m}\right)^{1/\gamma}.
\]
最后一个极限小于 $1$ 当且仅当
$\lambda m>L/2$，即 \eqref{eq:H-critical-threshold}。
\end{proof}


